%!TEX program = xelatex
\documentclass[a4paper,12pt]{article}
\usepackage{fontspec}\defaultfontfeatures{Ligatures=TeX}
\usepackage[a4paper,vmargin={3cm,3.5cm},hmargin={3cm,3cm}]{geometry}
%-----------------------------------------------------------------------------%
\usepackage{settings}
%-----------------------------------------------------------------------------%
%%% Title %%%
\title{G6411 Recitation Notes: Maximum Likelihood}
\author{Jesse Chieh Chen\\\href{mailto:cc5229@columbia.edu}{cc5229@columbia.edu}}
\date{\today}
%-----------------------------------------------------------------------------%

\begin{document}

\maketitle

\section{Likelihood and Maximum Likelihood}

Let $\{Y_i\}_{i=1}^n$ be \gls{iid},
and let $f(y\given\theta)$ be a proposed density,
where $\theta\in\Theta\subseteq\reals^k$.
For discrete outcomes, replace the density by a probability mass function.
The model is \emph{correctly specified} if, for some $\theta_0\in\Theta$,
$f(y\given\theta_0)$ is the true density of $Y_i$.
The true parameter $\theta_0$ is fixed but unknown.

\begin{definition}[Likelihood Function]
	The likelihood and average log-likelihood are
	\begin{align}
		L_n(\theta)
		&= \prod_{i=1}^nf(Y_i\given\theta) \\
		\shortintertext{and}
		\ell_n(\theta)
		&= \frac1n \log L_n(\theta)
		=\frac1n \sum_{i=1}^n\log f(Y_i\given\theta).
	\end{align}
\end{definition}

\begin{remark}[Density v.s.\ Likelihood Functions]
	In a density, the parameter is fixed and the outcome varies.
	In a likelihood, the observed data are fixed and we compare different parameter values.
	The likelihood is not a probability density for $\theta$,
	and its integral over $\Theta$ need not equal one.
\end{remark}

\begin{definition}[Likelihood and Maximum Likelihood]
	\Gls{mle} chooses a parameter that maximizes the likelihood:
	\begin{align}\label{eq-mle}
		\widehat\theta\in\arg\max_{\theta\in\Theta}L_n(\theta)
		=\arg\max_{\theta\in\Theta}\ell_n(\theta).
	\end{align}
\end{definition}

\section{Properties of Maximum Likelihood}

\subsection{Invariance}

\begin{proposition}[Invariance of \gls{mle}]
	Let $g$ be a one-to-one function of $\theta$.
	Define the induced likelihood for $\eta=g(\theta)$ by $L_n^\eta(\eta)\coloneqq L_n(g\inv(\eta))$.
	If $\widehat\theta$ maximizes $L_n(\theta)$,
	then $\widehat\eta\coloneqq g(\widehat\theta)$ maximizes $L_n^\eta(\eta)$.
\end{proposition}

\begin{remark}[Implication of Invariance]
	If $\widehat \sigma^2$ is the maximum likelihood estimate of a variance $\sigma^2$,
	then $\sqrt{\widehat \sigma^2}$ is the maximum likelihood estimate of the standard deviation $\sigma$.
	In general, invariance does not preserve unbiasedness.
	On the other hand, consistency is preserved if $g$ is continuous.
\end{remark}

\begin{remark}[No Jacobian for a Parameter Transformation]
	Reparameterizing the likelihood only relabels its argument.
	There is no Jacobian because we are not transforming a random variable with a density over $\Theta$.
	Transforming the \emph{data} is a different operation;
	that change-of-density formula is given in \autoref{app-density}.
\end{remark}

\begin{definition}[Population Objective]\label{def-population-objective}
	The population objective is the expected log-likelihood at each parameter value:
	\begin{align}
		Q(\theta)
		&\coloneqq \E[\log f(Y_i\given\theta)\given\theta_0]
		= \int_{\mathcal{Y}} \log f(y\given\theta)f(y\given\theta_0)\diff y.
	\end{align}
\end{definition}

\begin{remark}
	For conciseness, we write the population expectation as
	\begin{align}
		\E_0[\log f(Y_i\given\theta)]
		\coloneqq \E[\log f(Y_i\given\theta)\given\theta_0].
	\end{align}
	That is, we use $\E_0$ to denote expectation under the \emph{true} density of $Y_i$.
\end{remark}

\subsection{Score and Information}

\begin{definition}[Score and Hessian]
	Let $\ell_i(\theta)=\log f(Y_i\given\theta)$.
	The individual score is a column vector, and the individual Hessian is a matrix:
	\begin{gather}
		s_i(\theta)
		\coloneqq \dot\ell_i(\theta)
		\coloneqq \frac{\partial}{\partial\theta} \ell_i(\theta)
		= f(Y_i\given\theta)\inv\frac{\partial}{\partial\theta} f(Y_i\given\theta)
		\in\reals^k
		\shortintertext{and}
		H_i(\theta)
		\coloneqq \ddot\ell_i(\theta)
		\coloneqq \frac{\partial^2}{\partial\theta\partial\theta\T} \ell_i(\theta)
		\in\reals^{k\times k}.
	\end{gather}
\end{definition}

\begin{remark}[On Notation]
	We use dots to denote derivatives w.r.t.\ parameter $\theta$.
\end{remark}

\begin{remark}[Solving for \gls{mle} using the Score]
	Typically, solving for the maximization problem \eqref{eq-mle}
	means solving the first-order condition, i.e., the score equation
	\begin{align}\label{eq-score-equation}
		\frac{\partial\ell_n(\widehat\theta)}{\partial\theta}
		=\frac1n \sum_{i=1}^ns_i(\widehat\theta)=0
	\end{align}
	provided that $\widehat\theta$ is in the interior of $\Theta$ and the objective is differentiable there.
\end{remark}

\begin{definition}[Information Matrix]\label{def-information}
	The \emph{information matrix} is defined as
	\begin{align}
		I(\theta_0)\coloneqq\E_0[s_i(\theta_0)s_i(\theta_0)\T].
	\end{align}
	It is also known as the \emph{Fisher information}.
\end{definition}

\begin{assumption}[Regularity Conditions for Likelihood Calculations]\label{ass-regular}
	Here are some standard regularity conditions for the calculations below:
	\begin{enumerate}
		\item The true parameter $\theta_0$ is in the interior of $\Theta$.
		\item The support of the density does not depend on $\theta$.
		\item The log density is twice continuously differentiable near $\theta_0$.
		\item The score has finite second moments.
		\item The information matrix is positive definite.
	\end{enumerate}
	We also require that differentiation and integration can be interchanged in the calculations below.
\end{assumption}

\begin{proposition}
	$\E_0 s_i(\theta_0)=0$.
\end{proposition}
\begin{proof}
	Consider the following:
	\begin{align*}
		\E_0 s_i(\theta_0)
		&= \int_{\mathcal{Y}} \left.\frac{\partial}{\partial\theta}\log f(y\given\theta)\right|_{\theta=\theta_0} f(y\given\theta_0)\diff y \\
		&= \int_{\mathcal{Y}} \left.\frac{\partial}{\partial\theta} f(y\given\theta)\right|_{\theta=\theta_0} \diff y
		= \left.\frac{\partial}{\partial\theta}\int_{\mathcal{Y}} f(y\given\theta)\diff y\right|_{\theta=\theta_0}
		= 0.
		\qedhere
	\end{align*}
	The penultimate equality
	first exchanges differentiation and integration,
	then uses the fact that the integral of a density is one, so its derivative is zero.
\end{proof}

\begin{corollary}
	$\var_0(s_i(\theta_0))=I(\theta_0)$.
\end{corollary}

\begin{question}[Interpreting Fisher Information]\label{question-1}
	Knowing that information matrix is the variance of the score,
	is it good or bad to have a \emph{large} information matrix?
\end{question}

\begin{theorem}[Information Equality]\label{thm-information-equality}
	$I(\theta_0)=-\E_0 H_i(\theta_0)$.
\end{theorem}
\begin{proof}
	Consider the following:
	\begin{align}
		\E_0 H_i(\theta_0)
		&= \int_{\mathcal{Y}} \left.\frac{\partial^2}{\partial\theta\partial\theta\T}\log f(y\given\theta)\right|_{\theta=\theta_0} f(y\given\theta_0)\diff y \\
		&= \int_{\mathcal{Y}} \left(\frac{\ddot f(y\given\theta_0)}{f(y\given\theta_0)} - \frac{\dot f(y\given\theta_0)\dot f(y\given\theta_0)\T}{f(y\given\theta_0)^2}\right) f(y\given\theta_0)\diff y \\
		&= \underbrace{\left.\frac{\partial^2}{\partial\theta\partial\theta\T}\int_{\mathcal{Y}} f(y\given\theta)\diff y\right|_{\theta=\theta_0}}_{=0}
		- \int_{\mathcal{Y}} \frac{\dot f(y\given\theta_0)\dot f(y\given\theta_0)\T}{f(y\given\theta_0)}\diff y
		= - I(\theta_0).
	\end{align}
	The penultimate equality
	first exchanges differentiation and integration,
	then uses the fact that the integral of a density is one, so its derivative is zero.
\end{proof}

\begin{question}[Interpreting Fisher Information]\label{question-2}
	Knowing that information matrix is the expected negative curvature of the log-likelihood,
	is it good or bad to have a \emph{large} information matrix?
	How does your answer compare to \autoref{question-1}?
\end{question}

\subsection{Consistency}

The key convergence condition for consistency is that the average log-likelihood
converges to its expectation \emph{uniformly} over $\Theta$:
\begin{align}\label{eq-uniform-convergence}
	\sup_{\theta\in\Theta}\left|\ell_n(\theta)-Q(\theta)\right|\pto 0.
\end{align}
\begin{proposition}[Consistency of \gls{mle}]\label{prop-mle-consistency}
	Suppose also that $\Theta$ is compact and $Q$ is continuous with unique maximizer $\theta_0$.
	Then, \eqref{eq-uniform-convergence} guarantees consistency.
\end{proposition}
\begin{proof}
	Note that $\ell_n(\widehat\theta)\ge\ell_n(\theta_0)$ by definition of $\widehat\theta$.
	Then,
	\begin{align*}
		0\le Q(\theta_0)-Q(\widehat\theta)
		\le |Q(\theta_0)-\ell_n(\theta_0)|
		+|\ell_n(\widehat\theta)-Q(\widehat\theta)| 
		\le 2\sup_{\theta\in\Theta}|\ell_n(\theta)-Q(\theta)|\pto 0.
	\end{align*}
	For any $\varepsilon>0$, compactness, continuity, and uniqueness imply a positive gap
	$\delta_\varepsilon=Q(\theta_0)-\sup_{\theta\in\Theta:\|\theta-\theta_0\|\ge\varepsilon}Q(\theta)>0$
	whenever the set is nonempty.
	Thus
	\begin{align*}
		\Pr(\|\widehat\theta-\theta_0\|\ge\varepsilon)
		&\le \Pr\big(Q(\theta_0)-Q(\widehat\theta)\ge\delta_\varepsilon\big)\to 0.
		\qedhere
	\end{align*}
\end{proof}

\begin{remark}
	Uniform convergence is usually established using a uniform law of large numbers.
	Common sufficient conditions are a compact parameter space, continuity of $\log f(y\given\theta)$ in $\theta$,
	and a bound $|\log f(Y_i\given\theta)|\le M(Y_i)$ for all $\theta\in\Theta$, with $\E_0M(Y_i)<\infty$.
	These conditions control how the sample objective varies across parameter values,
	so convergence holds uniformly rather than only at each fixed $\theta$.
\end{remark}

\subsection{Asymptotic Normality}\label{sec-mle-asymptotic-normality}

Suppose that the \gls{mle} $\widehat\theta$ is consistent.
Also assume that the sample average Hessian converges uniformly near $\theta_0$ to its continuous expectation.
Then, by the integral form of the mean value expansion, we have
\begin{gather}
	\dot\ell_n(\widehat\theta) - \dot\ell_n(\theta_0) = \widetilde H_n(\widehat\theta-\theta_0),
	\shortintertext{where}
	\widetilde H_n
	=\int_0^1\ddot\ell_n\big(\theta_0+t(\widehat\theta-\theta_0)\big)\diff t
	\pto\E_0 H_i(\theta_0).
\end{gather}
Note that $\dot\ell_n(\widehat\theta)=0$ for an interior maximizer.
Consistency and interiority of $\theta_0$ ensure that $\widehat\theta$ is interior with probability approaching one.
Then, by the \gls{clt} and the Hessian limit above, we have
\begin{align}
	\sqrt{n}(\widehat\theta-\theta_0)
	&= - \widetilde H_n\inv \frac{1}{\sqrt{n}}\sum_{i=1}^n s_i(\theta_0) \\
	&\dto
	- \E_0[H_i(\theta_0)]^{-1} \normaldistribution(0, I(\theta_0))
	= \normaldistribution(0, V). \label{eq-mle-normality}
	\shortintertext{where}
	V &=  \E_0[H_i(\theta_0)]^{-1} I(\theta_0) \E_0[H_i(\theta_0)]^{-1} = I(\theta_0)\inv\label{eq-mle-asymptotic-variance}
\end{align}
where the simplification of asymptotic variance is by \nameref{thm-information-equality}.

\begin{remark}
	Answers to \autoref{question-1} and \autoref{question-2} are consistent.
	Note that in the asymptotic variance formula, the variance of $s_i(\theta_0)$, i.e., $I(\theta_0)$, is in the numerator,
	whereas the curvature $\E_0[H_i(\theta_0)]$ is in the denominator.
	Efficiency of \gls{mle} is a trade-off between the two quantities.
	In the end, we want a \emph{large} Fisher information matrix, hence the name ``information.''
	The intuition of \eqref{eq-mle-normality} will carry over to the \gls{qmle} later.
\end{remark}

% \begin{remark}[Asymptotic versus Finite Sample Variance]
% 	Equation \eqref{eq-mle-normality} suggests the large sample covariance approximation
% 	$I(\theta_0)\inv/n$ for $\widehat\theta$.
% 	It is not generally an exact finite sample covariance formula,
% 	and maximum likelihood estimates need not be unbiased.
% \end{remark}

\subsection{The Cram\'er-Rao Bound}

\begin{theorem}[Cram\'er-Rao Bound]\label{thm-cramer-rao}
	In a correctly specified regular model,
	let $T$ have finite variance and satisfy $\E_\theta T=\theta$
	for every $\theta$ in a neighborhood of $\theta_0$.
	Assume differentiation can also be passed through the expectation of $T$.
	Then
	\begin{align}
		\var_0(T)\succeq[nI(\theta_0)]\inv.
	\end{align}
\end{theorem}
\begin{proof}
	Since $\E_\theta T=\theta$ for all $\theta$ near $\theta_0$,
	differentiating with respect to $\theta$ gives
	\begin{gather}\label{eq-cramer-rao-derivative}
		\int_{\mathcal{Y}^n} T
		\underbrace{\left[\frac{\dot f(y_1,...,y_n\given\theta)\T}{f(y_1,...,y_n\given\theta)}\right]}_{S_n\T}
		f(y_1,...,y_n\given\theta) \diff y_1...\diff y_n = I_k,
	\end{gather}
	evaluating at $\theta_0$ gives $\E_0 T S_n\T = I_k$,
	where $S_n=\sum_i s_i(\theta_0)$ is the sum of scores.
	Let $\var_0(S_n)=J_n=nI(\theta_0)$.
	Hence, we have the following:
	\begin{align*}
		0
		&\preceq \var_0(T-J_n\inv S_n)
		= \var_0(T)-J_n\inv.
		\qedhere
	\end{align*}
\end{proof}

\begin{remark}
	The Cram\'er-Rao bound gives a lower limit on the variance of an unbiased estimator
	in each direction. An unbiased estimator that attains the bound has the smallest possible
	variance within the class covered by the theorem, although the bound need not be attainable.
	Unlike Gauss-Markov, the comparison is not restricted to estimators linear in the data.
	Instead, it applies to unbiased estimators satisfying the relevant regularity conditions
	in a correctly specified model.
	Roughly, these conditions require the density to vary smoothly with the parameter
	and allow differentiation to pass through the expectations used in the proof.
\end{remark}

\section{The Normal Linear Model}

\subsection{Likelihood and Maximization}

Consider
\begin{align}
	Y_i=X_i\T\beta_0+e_i,
	\qquad e_i\given X_i\sim\normaldistribution(0,\sigma_0^2),
	\qquad \sigma_0^2>0.
\end{align}
The observations are \gls{iid}, $X_i\in\reals^k$,
and the sample matrix $X$ has rows $X_i\T$ and full column rank.
Derivatives below are with respect to the variance $\sigma^2$,
not the standard deviation $\sigma$.
Here $\beta\in\reals^k$, so the full parameter $\theta=(\beta\T,\sigma^2)\T$ has dimension $k+1$.
For a candidate $\theta$,
\begin{align}
	f(y_i\given x_i,\beta,\sigma^2)
	&=(2\pi \sigma^2)^{-1/2}\exp\left\{-\frac{(y_i-x_i\T\beta)^2}{2\sigma^2}\right\},\\
	\ell_n(\beta,\sigma^2)
	&=-\frac12\log(2\pi)-\frac12\log \sigma^2
	-\frac1{2n\sigma^2}(Y-X\beta)\T(Y-X\beta).
\end{align}
The two score equations are
\begin{align}
	\frac{\partial\ell_n}{\partial\beta}
	&=\frac{X\T(Y-X\beta)}{n\sigma^2}=0,\\
	\frac{\partial\ell_n}{\partial \sigma^2}
	&=-\frac1{2\sigma^2}+\frac{(Y-X\beta)\T(Y-X\beta)}{2n\sigma^4}=0.
\end{align}
For any fixed $\sigma^2>0$, maximizing the log-likelihood over $\beta$ is equivalent to minimizing
the sum of squared residuals. Hence
\begin{align}\label{eq-normal-mle}
	\widehat\beta=(X\T X)\inv X\T Y,
	\qquad
	\widehat\sigma^2=\frac{\widehat e\T\widehat e}{n},
	\qquad \widehat e=Y-X\widehat\beta.
\end{align}
After substituting $\widehat\beta$, the derivative with respect to $\sigma^2$ has the sign of
$\widehat e\T\widehat e-n\sigma^2$.
Thus $\widehat\sigma^2$ is the global maximizer over $\sigma^2>0$ whenever $\widehat e\T\widehat e>0$.
Assume $n>k$ and exclude a perfect fit; under the stated normal model a perfect fit has probability zero.

\begin{remark}[Why Condition on the Regressors?]
	We model the distribution of $Y_i$ given $X_i$, rather than their joint distribution.
	This lets us estimate $(\beta,\sigma^2)$ without specifying how the regressors are distributed.
	When the distribution of $X_i$ does not depend on $(\beta,\sigma^2)$,
	its contribution to the joint likelihood is constant in these parameters,
	so maximizing the conditional likelihood gives the same estimates.
\end{remark}

\begin{remark}[Why Divide by $n$, Not $n-k$?]
	The divisor $n$ comes from maximizing the likelihood.
	Since $\widehat e=M e$ with $M=I_n-X(X\T X)\inv X\T$,
	\begin{align*}
		\E(\widehat e\T\widehat e\given X)
		=\sigma_0^2\operatorname{tr}(M)=(n-k)\sigma_0^2.
	\end{align*}
	Consequently, $\widehat\sigma^2$ is biased downward:
	$\E(\widehat\sigma^2\given X)=(n-k)\sigma_0^2/n$.
	The estimator $\widehat e\T\widehat e/(n-k)$ is unbiased but is not the maximum likelihood estimate.
\end{remark}

\subsection{Hessian and Information}

Writing $e_i(\beta)=Y_i-X_i\T\beta$, the individual score and Hessian are
\begin{align}
	s_i(\beta,\sigma^2)
	&=\begin{pmatrix}
		X_i e_i(\beta)/\sigma^2\\
		-1/(2\sigma^2)+e_i(\beta)^2/(2\sigma^4)
	\end{pmatrix},\\
	H_i(\beta,\sigma^2)
	&=\begin{pmatrix}
		-X_iX_i\T/\sigma^2 & -X_i e_i(\beta)/\sigma^4\\
		-e_i(\beta)X_i\T/\sigma^4 & 1/(2\sigma^4)-e_i(\beta)^2/\sigma^6
	\end{pmatrix}.
\end{align}
Let $Q_{xx}=\E(X_iX_i\T)$ be finite and positive definite.
At the truth, $\E(e_i\given X_i)=0$ and $\E(e_i^2\given X_i)=\sigma_0^2$, so
\begin{align}\label{eq-normal-information}
	I(\theta_0)=-\E_0H_i(\theta_0)
	=\begin{pmatrix}
		Q_{xx}/\sigma_0^2&0\\
		0&1/(2\sigma_0^4)
	\end{pmatrix}.
\end{align}
We can also calculate the score covariance directly.
Normality gives $\E(e_i^3\given X_i)=0$ and $\E(e_i^4\given X_i)=3\sigma_0^4$.
Therefore
\begin{align*}
	\E_0(s_{\beta i}s_{\beta i}\T)&=Q_{xx}/\sigma_0^2,\\
	\E_0(s_{\beta i}s_{\sigma^2,i})&=0,\\
	\E_0(s_{\sigma^2,i}^2)&=\frac{\E_0[(e_i^2-\sigma_0^2)^2]}{4\sigma_0^8}
	=\frac{3\sigma_0^4-\sigma_0^4}{4\sigma_0^8}=\frac1{2\sigma_0^4}.
\end{align*}
This verifies the information equality for both parameters.
Inverting \eqref{eq-normal-information} gives
\begin{align}
	\sqrt n\begin{pmatrix}\widehat\beta-\beta_0\\\widehat\sigma^2-\sigma_0^2\end{pmatrix}
	\dto\normaldistribution\left(0,
	\begin{pmatrix}\sigma_0^2Q_{xx}\inv&0\\0&2\sigma_0^4\end{pmatrix}\right).
\end{align}

\begin{remark}[Conditioning and the Asymptotic Distribution]
	Although we construct the likelihood conditional on the regressors,
	the asymptotic calculation here treats $(Y_i,X_i)$ as \gls{iid} random observations.
	Thus, the expectations defining the information matrix average over both $Y_i$ and $X_i$.
	In particular, the conditional information for $\beta$ is $X_iX_i\T/\sigma_0^2$,
	while its expectation is $Q_{xx}/\sigma_0^2$, where $Q_{xx}=\E(X_iX_i\T)$.
\end{remark}

\begin{remark}
	Conditional on $X$, the sample information for $\beta$ is $X\T X/\sigma_0^2$.
	Here the inverse information is also the \emph{exact} conditional covariance:
	\begin{align}
		\widehat\beta\given X
		\sim\normaldistribution\big(\beta_0,\sigma_0^2(X\T X)\inv\big).
	\end{align}
	Thus $\widehat\beta$ attains the conditional Cram\'er-Rao bound in this model,
	even when $\sigma_0^2$ is unknown, because the information matrix is block diagonal.
\end{remark}

\section{Three Routes to \glsentryshort{ols}}

The same estimator can be motivated by three different problems:
\begin{enumerate}
	\item \emph{Loss minimization.}
	Choose the linear predictor that minimizes sample squared error:
	$\min_b\sum_i(Y_i-X_i\T b)^2$.
	This is an algebraic optimization problem, requiring no distributional assumptions.
	\item \emph{Method of moments.}
	Start with the population restriction $\E[X_i(Y_i-X_i\T\beta_0)]=0$
	and impose its sample analogue $n\inv\sum_iX_i(Y_i-X_i\T\widehat\beta)=0$.
	This only specifies particular moments, not a full distribution.
	\item \emph{Maximum likelihood.}
	Assume a normal conditional distribution with linear mean and constant variance,
	and maximize its likelihood as in \eqref{eq-normal-mle}.
\end{enumerate}
Full column rank makes each sample problem give $\widehat\beta=(X\T X)\inv X\T Y$.
The formula is identical, but the assumptions used to interpret it are not.
In particular, normality is not needed for the usual consistency of \gls{ols}.

\section{Quasi-Maximum Likelihood}

\subsection{What If the Density Is Wrong?}

Let $g(y\given x)$ denote the \emph{true conditional density} of $Y_i$ given $X_i=x$,
and let $f(y\given x,\theta)$ be a \emph{working conditional density} chosen for estimation.
Under correct specification, $g(y\given x)=f(y\given x,\theta_0)$ for some $\theta_0$.
% Under misspecification, no such parameter need exist.
We maximize the working likelihood $f$, but evaluate the estimator's properties under the true distribution $g$.
Write $\E_g$ for expectation under the true distribution of $(Y_i,X_i)$.
The estimator is
\begin{align}
	\widehat\theta\in\arg\max_{\theta\in\Theta}
	\frac1n\sum_{i=1}^n\log f(Y_i\given X_i,\theta).
\end{align}
This is \emph{\gls{qmle}}.
Its population target is the \emph{pseudo-true parameter}
\begin{align}
	\theta_*
	\in \arg\max_{\theta\in\Theta} \E_g[\log f(Y_i\given X_i,\theta)].
\end{align}
Under identification and suitable uniform convergence, $\widehat\theta\pto\theta_*$.
There is no general guarantee that $\theta_*$ equals a parameter of interest.
That must be checked for the particular model and the particular misspecification.

\subsection{A Normal Working Likelihood under Heteroskedasticity}

Suppose the true model satisfies
\begin{align}\label{eq-heteroskedastic-model}
	Y_i=X_i\T\beta_0+e_i,
	\qquad \E_g(e_i\given X_i)=0,
	\qquad \var_g(e_i\given X_i)=\sigma^2(X_i),
\end{align}
but our working density $f$ is the same homoskedastic normal density as before.
Here $g(y\given x)$ is the true conditional density with mean $x\T\beta_0$
and variance $\sigma^2(x)$; it need not be normal.
Assume finite second moments of $Y_i$ and $X_i$ and positive definite $Q_{xx}$.
The sample solution remains \eqref{eq-normal-mle}.
To find its population target, write
\begin{align}
	\E_g[(Y_i-X_i\T\beta)^2]
	=\E_g(e_i^2)+(\beta-\beta_0)\T Q_{xx}(\beta-\beta_0).
\end{align}
The cross term vanishes because $\E_g(X_ie_i)=0$.
Hence expected squared error is uniquely minimized at $\beta_0$, and the pseudo-true parameters are
\begin{align}
	\beta_*=\beta_0,
	\qquad \sigma_*^2=\E_g(e_i^2)=\E_g[\sigma^2(X_i)]>0.
\end{align}
Thus \gls{ols} still consistently estimates $\beta_0$,
while $\widehat\sigma^2$ consistently estimates the \emph{average} error variance,
not the observation-specific conditional variance.
Neither normality nor homoskedasticity is needed for these consistency conclusions.

\subsection{Why the Information Equality Can Fail}

At an interior pseudo-true parameter, define
\begin{align}
	A=-\E_gH_i(\theta_*),
	\qquad B=\E_g[s_i(\theta_*)s_i(\theta_*)\T].
\end{align}
Under the differentiability and moment conditions needed for the score expansion,
$\E_gs_i(\theta_*)=0$.
The score and Hessian are derivatives of $\log f(Y_i\given X_i,\theta)$,
but these expectations are taken under $g$.
Recall the proof of \nameref{thm-information-equality}:
the equality obtained by integrating against $f$ need not hold under $g$,
so there is no general reason for $A=B$.
With nonsingular $A$ and the usual consistency, Hessian convergence, and score \gls{clt} conditions,
\begin{align}\label{eq-sandwich}
	\sqrt n(\widehat\theta-\theta_*)
	\dto\normaldistribution(0,A\inv BA\inv).
\end{align}
This is the \emph{sandwich} covariance:
curvature on the outside and score variability in the middle.

\begin{remark}
	Recall the score expansion in \autoref{sec-mle-asymptotic-normality}:
	the inverse Hessian appears on both sides of the score covariance.
	The same argument applies here, but without the information equality we cannot simplify $A\inv BA\inv$ to $A\inv$.
	Under correct specification, $A=B=I(\theta_0)$ and it reduces to \eqref{eq-mle-asymptotic-variance}.
\end{remark}

For the normal working likelihood in \eqref{eq-heteroskedastic-model},
let $g_i=X_ie_i$ and suppose $\Omega=\E_g[g_i g_i\T]$ is finite.
The matrices are
\begin{align}
	A &= 
	\begin{bmatrix*}
		A_{\beta\beta} & 0\\
		0 & A_{\sigma^2\sigma^2}
	\end{bmatrix*}
	=
	\begin{bmatrix*}
		Q_{xx}/\sigma_*^2 & 0\\
		0 & 1/(2\sigma_*^4)
	\end{bmatrix*} \\
	\shortintertext{and}
	B &=
	\begin{bmatrix*}
		B_{\beta\beta} & B_{\beta,\sigma^2}\\
		B_{\beta,\sigma^2}\T & B_{\sigma^2\sigma^2}
	\end{bmatrix*}
	=
	\begin{bmatrix*}
		\Omega/\sigma_*^4 & \E_g[X_ie_i^3]/(2\sigma_*^6)\\
		\E_g[e_i^3X_i\T]/(2\sigma_*^6) & \var_g(e_i^2)/(4\sigma_*^8)
	\end{bmatrix*}.
\end{align}
Consequently, the slope and intercept coefficients have limiting covariance
\begin{align}
	A_{\beta\beta}\inv B_{\beta\beta}A_{\beta\beta}\inv
	=Q_{xx}\inv\Omega Q_{xx}\inv.
\end{align}
The inverse-Hessian formula alone would instead give $\sigma_*^2Q_{xx}\inv$,
which is generally incorrect under heteroskedasticity.
The corresponding heteroskedasticity-robust covariance estimator for $\widehat\beta$ is
\begin{align}
	\widehat{\var}(\widehat\beta)
	=(X\T X)\inv\left(\sum_{i=1}^nX_iX_i\T\widehat e_i^2\right)(X\T X)\inv.
\end{align}
Its consistency requires the usual additional moment conditions for this sample sandwich.
The joint asymptotic distribution including $\widehat\sigma^2$ also requires moments of $e_i^2$,
in particular a finite fourth moment of $e_i$ for its score variance.

\begin{remark}[Why Use a Working Likelihood?]
	A convenient likelihood can provide a tractable objective and useful estimating equations
	without requiring its entire distributional specification to be correct.
	% The normal example works because the coefficient score is proportional to $X_ie_i$,
	% whose expectation is zero under the conditional mean restriction.
	Other working likelihoods can similarly target a correctly specified conditional mean.
	Consistency must still be justified, and standard errors must allow for the misspecification.
	There is no general efficiency guarantee from choosing a convenient working likelihood.
\end{remark}

\vfill
\section*{Acronyms}
\printglossaries

\clearpage
\appendix

\section{Change of Density}\label{app-density}

\begin{theorem}[Change of Density]
	Let $h$ be a one-to-one differentiable function with differentiable inverse, and let
	$Z\coloneqq h(Y)$.
	Let $f_{Y\given X}(y\given x,\theta)$ and $f_{Z\given X}(z\given x,\theta)$ denote the conditional densities
	of $Y$ given $X$ and $Z$ given $X$, respectively.
	Then the two densities are related by the change-of-density formula
	\begin{align}\label{eq-change-of-density}
		f_{Z\given X}(z\given x,\theta)
		=f_{Y\given X}(h\inv(z)\given x,\theta)
		\left|\frac{\diff}{\diff z}h\inv(z)\right|.
	\end{align}
\end{theorem}

\begin{remark}[Change-of-Density via \glsentryshort{cdf}]
	Let $g$ be an arbitrary function of $Y$.
	Integrals of $g$ can be viewed as Riemann-Stieltjes integrals with respect to the \glspl{cdf}.
	Since $Z=h(Y)$, we have
	\begin{align}
		\int_{\mathcal{Y}} g(y)\diff F_{Y\given X}(y\given x)
		&=
		\int_{\mathcal{Z}} g(h\inv(z))\diff F_{Z\given X}(z\given x).
	\end{align}
	Now suppose that $F_{Y\given X}$ is differentiable with derivative $f_{Y\given X}$,
	and write the integral with respect to the density instead of the \gls{cdf}.
	The \gls{cdf} of $Z$ given $X$ is
	\begin{align}
		F_{Z\given X}(z\given x)
		= \Pr(Z\le z\given X=x)
		= \Pr(h(Y)\le z\given X=x).
	\end{align}
	We have two cases:
	\begin{align}
		F_{Z\given X}(z\given x)
		&=
		\begin{cases}
			F_{Y\given X}(h\inv(z)\given x)
			& \text{if $h$ is increasing,} \\
			1-F_{Y\given X}(h\inv(z)\given x)
			& \text{if $h$ is decreasing,}
		\end{cases} \\
		\implies
		\diff F_{Z\given X}(z\given x)
		&=
		\begin{cases}
			f_{Y\given X}(h\inv(z)\given x)
			\displaystyle\frac{\diff}{\diff z}h\inv(z)\diff z
			& \text{if $h$ is increasing,} \\
			-f_{Y\given X}(h\inv(z)\given x)
			\displaystyle\frac{\diff}{\diff z}h\inv(z)\diff z
			& \text{if $h$ is decreasing.}
		\end{cases}
	\end{align}
	Since the derivative of $h\inv$ is positive in the first case and negative
	in the second, both cases can be written as
	\begin{align}
		\diff F_{Z\given X}(z\given x)
		=
		f_{Y\given X}(h\inv(z)\given x)
		\left|
			\frac{\diff}{\diff z}h\inv(z)
		\right|
		\diff z,
	\end{align}
	which yields \eqref{eq-change-of-density}.
\end{remark}

\begin{remark}[Vector-Valued Transformations]
	For a vector-valued transformation, replace the absolute derivative by the absolute determinant
	of the inverse transformation's Jacobian matrix.
	Let $Y\in\reals^{k}$ and $Z\coloneqq h(Y)\in\reals^{k}$, and let
	$J_{h\inv}(z)\in\reals^{k\times k}$ denote the Jacobian matrix of $h\inv$ at $z$.
	Then
	\begin{align}
		f_{Z\given X}(z\given x,\theta)
		=f_{Y\given X}(h\inv(z)\given x,\theta)
		\left|\det J_{h\inv}(z)\right|.
	\end{align}
\end{remark}

\end{document}
